%%
%% Blank IEEEtran journal article skeleton
%% Stripped of documentation/instructional text — structure only.
%%
%% \documentclass[journal,onecolumn]{IEEEtran} is for IEEE journals/
%% transactions using a single-column layout (e.g. for review copies,
%% or venues that request onecolumn). Drop "onecolumn" for the standard
%% two-column journal layout:
%%   \documentclass[journal]{IEEEtran}
%% For a conference paper, drop the "journal" option:
%%   \documentclass{IEEEtran}
%% For IEEE Computer Society conferences/journals, add "compsoc":
%%   \documentclass[journal,compsoc]{IEEEtran}
%%
%% Refer to your target journal's author instructions for the exact
%% options it requires.
\documentclass[journal,onecolumn]{IEEEtran}

%% Common optional packages (uncomment as needed):
%%\usepackage{amsmath}          %% math environments/symbols
%%\interdisplaylinepenalty=2500 %% restore page breaks in multiline eqns after amsmath
%%\usepackage{cite}             %% auto-sorts/compresses \cite{} numbers, e.g. [1]-[3]
%%\usepackage{graphicx}         %% required for \includegraphics
%%\usepackage{algorithmic}      %% algorithm environment (do NOT use algorithm/algorithm2e floats)
%%\usepackage{array}            %% improved tabular/array environments
%%\usepackage[caption=false,font=footnotesize]{subfig} %% subfigures — use caption=false
%%                                                       %% to preserve IEEE-style captions
%%\usepackage{url}              %% \url{} for breakable hyperlinks

%% *** Do not adjust margins/column widths, and do not use font-altering
%% packages (e.g. pslatex) unless the venue explicitly requires it. ***

%% Correct hyphenation of specific words as needed:
%%\hyphenation{op-tical net-works semi-conduc-tor}

\begin{document}

%% Paper title. Capitalize per title-case conventions; no math/symbols in
%% the title. Use \\ for a manual linebreak if needed.
\title{}

%% ---------------------------------------------------------------
%% Authors and IEEE memberships.
%% Use ~ (nonbreaking space) between name parts so a name doesn't
%% break across lines. Use \thanks{} for affiliation/funding/manuscript
%% date footnotes — a separate \thanks{} per paragraph of footnote text.
%% Trailing % after \IEEEmembership/\thanks lines prevents stray spaces.
%% ---------------------------------------------------------------
\author{%
  ~\IEEEmembership{}%
\thanks{}%
\thanks{Manuscript received ; revised .}}

%% Running headers: left = journal name/vol/no/date, right = short author/title.
%% Only appears on odd pages in twoside mode. Omit author names here for
%% blind peer review.
\markboth{}%
{}

%% Publisher ID mark on the page (optional):
%%\IEEEpubid{}
%% If used, call \IEEEpubidadjcol in the second column so text clears it.

%% Special paper notice, e.g. "(Invited Paper)" (optional):
%%\IEEEspecialpapernotice{}

\maketitle

%% No math, special symbols, or citations in the abstract/keywords.
\begin{abstract}

\end{abstract}

%% Keywords are not normally used for peer-review-mode papers.
\begin{IEEEkeywords}

\end{IEEEkeywords}

%% Extra cover-page info for peer-review mode (optional):
%%\ifCLASSOPTIONpeerreview
%%\begin{center} \bfseries EDICS Category: \end{center}
%%\fi

%% Inserts a page break + second title page in peer-review mode;
%% ignored otherwise. Keep this call regardless of mode.
\IEEEpeerreviewmaketitle

\section{Introduction}
%% Optional IEEE drop-cap opener for the first paragraph:
%%\IEEEPARstart{T}{his} paper ...

\subsection{}

\subsubsection{}

%% needed in the second column of the first page if using \IEEEpubid:
%%\IEEEpubidadjcol

%% ---------------------------------------------------------------
%% Figure example. \label must come AFTER \caption. IEEE convention:
%% floats are typically placed at the top of the column/page only.
%% ---------------------------------------------------------------
\begin{figure}[!t]
\centering
\includegraphics[width=2.5in]{}
\caption{}
\label{fig:1}
\end{figure}

%% Double-column (spanning) figure with two subfigures — requires subfig.sty:
%%\begin{figure*}[!t]
%%\centering
%%\subfloat[]{\includegraphics[width=2.5in]{}%
%%\label{fig:1a}}
%%\hfil
%%\subfloat[]{\includegraphics[width=2.5in]{}%
%%\label{fig:1b}}
%%\caption{}
%%\label{fig:2}
%%\end{figure*}

%% ---------------------------------------------------------------
%% Table example. IEEE convention: \caption comes BEFORE the table,
%% and \label comes after \caption. Table text defaults to \footnotesize.
%% ---------------------------------------------------------------
\begin{table}[!t]
\renewcommand{\arraystretch}{1.3}
\caption{}
\label{tab:1}
\centering
\begin{tabular}{|c||c|}
\hline
 & \\
\hline
 & \\
\hline
\end{tabular}
\end{table}

\section{}

\section{Conclusion}

%% ---------------------------------------------------------------
%% Appendices.
%% Single appendix: \appendix[Title] or \appendix (no heading).
%% Multiple appendices: \appendices, then \section{} for each —
%% do NOT use \section after \appendix/\appendices, only \section*
%% if additional unnumbered headings are needed within.
%% ---------------------------------------------------------------
\appendices
\section{}

\section{}

%% Acknowledgment uses \section* (not \section):
\section*{Acknowledgment}

%% ---------------------------------------------------------------
%% Bibliography.
%% Option A — BibTeX-generated:
%%\bibliographystyle{IEEEtran}
%%\bibliography{}
%% Option B — manually authored thebibliography (2nd \begin arg =
%% number of entries, reserves space for the numbering box):
%% ---------------------------------------------------------------
\begin{thebibliography}{1}

\bibitem{}

\end{thebibliography}

%% ---------------------------------------------------------------
%% Biography section (optional, common in IEEE Transactions).
%% With photo: pass \includegraphics as the optional argument, wrapped
%% in an extra set of braces to avoid parser confusion:
%%\begin{IEEEbiography}[{\includegraphics[width=1in,height=1.25in,clip,keepaspectratio]{}}]{Author Name}
%% Without a photo, reserve blank space:
%%\begin{IEEEbiography}{Author Name}
%%
%%\end{IEEEbiography}
%% If no photo at all will ever be included, use the nophoto variant:
%%\begin{IEEEbiographynophoto}{Author Name}
%%
%%\end{IEEEbiographynophoto}
%% ---------------------------------------------------------------

\end{document}
%% End of blank IEEEtran journal (one-column) skeleton.
